\chapter{Bilangan Kompleks}\label{ch:b1:complex}

Bilangan kompleks sudah ditemui pada jilid sekolah menengah sebagai
perkakas hitung untuk persamaan kuadrat. Bab ini memperlakukannya
sebagai objek pusat: bentuk eksponensial beserta akibatnya (de Moivre,
akar ke-$n$, \omterm{def:b1:complex:unity}{akar satuan}), penerjemahan sistematis antara $\C$ dan
geometri bidang, serta pemakaian $\eu^{\iu\theta}$ sebagai mesin
pembukti kesamaan trigonometri.

\section{Lapangan $\C$, modulus dan konjugat}

\begin{definition}[Lapangan kompleks]\label{def:b1:complex:field}
$\C = \{a + \iu b : a, b \in \R\}$ dengan penjumlahan biasa dan
perkalian yang ditentukan oleh $\iu^2 = -1$. Setiap $z = a + \iu
b$ yang tak nol mempunyai invers: $z^{-1} = \frac{a - \iu b}{a^2 + b^2}$ --- dalam
bahasa \cref{ch:b1:structures}, $\C$ adalah sebuah lapangan. Kita tulis $a =
\Re(z)$, $b = \Im(z)$, $\conj{z} = a - \iu b$ (yang disebut
\emph{konjugat}\index{konjugat}) dan $\abs{z} = \sqrt{a^2 + b^2}$
(yang disebut \emph{modulus}\index{modulus}).
\end{definition}

\begin{proposition}[Kaidah untuk konjugat dan modulus]\label{prop:b1:complex:rules}
Untuk $z, w \in \C$:
\begin{enumerate}
  \item $\conj{z + w} = \conj z + \conj w$, $\conj{zw} = \conj z\,
        \conj w$, $\conj{\conj z} = z$;
  \item $z \conj z = \abs z^2$; $\;\Re(z) = \frac{z + \conj z}{2}$,
        $\Im(z) = \frac{z - \conj z}{2\iu}$;
  \item $\abs{zw} = \abs z\, \abs w$, dan $\abs{z^{-1}} =
        \abs{z}^{-1}$ untuk $z \neq 0$;
  \item (ketaksamaan segitiga\index{ketaksamaan segitiga}) $\abs{z + w}
        \leq \abs z + \abs w$, dengan kesamaan jika dan hanya jika $z$ dan
        $w$ terletak pada satu sinar garis yang sama dari $0$ ($w = \lambda z$ atau
        $z = \lambda w$ dengan $\lambda \geq 0$);
  \item (ketaksamaan segitiga terbalik) $\bigl|\abs z - \abs w\bigr|
        \leq \abs{z - w}$.
\end{enumerate}
\end{proposition}

\begin{proof}
(1) dan (2) adalah perhitungan langsung pada bagian real dan bagian imajiner. (3):
$\abs{zw}^2 = zw\,\conj{zw} = z\conj z\, w \conj w = \abs z^2 \abs
w^2$, lalu tarik akar kuadratnya; terapkan pada $z \cdot z^{-1} = 1$ untuk
inversnya.

(4) Kedua ruasnya tak negatif, jadi bandingkan kuadratnya:
\[
\abs{z+w}^2 = (z+w)(\conj z + \conj w)
= \abs z^2 + \abs w^2 + 2\,\Re(z \conj w),
\]
dan $\Re(z\conj w) \leq \abs{z \conj w} = \abs z \abs w$ memberikan
$\abs{z+w}^2 \leq (\abs z + \abs w)^2$. Kesamaan memaksa $\Re(z \conj
w) = \abs{z \conj w}$, yaitu $z \conj w \in \R_+$; jika $w \neq 0$ hal ini
memberikan $z = \frac{z\conj w}{\abs w^2}\, w = \lambda w$ dengan $\lambda
\geq 0$ (dan kasus $w = 0$ bersifat sepele).

(5) $\abs z = \abs{(z - w) + w} \leq \abs{z-w} + \abs w$ memberikan $\abs z
- \abs w \leq \abs{z - w}$; tukarkan $z$ dan $w$ untuk tanda yang satunya.
\end{proof}

\begin{example}[Latihan lengkap modulus dan argumen]\label{ex:b1:complex:workout}
Nyatakan $w = \dfrac{3 + 4\iu}{1 - 2\iu}$ dalam bentuk aljabar lalu hitung
modulusnya dengan dua cara. Dengan mengalikan \omterm{def:b1:complex:field}{konjugat}
penyebutnya:
\[
w = \frac{(3 + 4\iu)(1 + 2\iu)}{(1 - 2\iu)(1 + 2\iu)}
= \frac{3 + 6\iu + 4\iu - 8}{1 + 4}
= \frac{-5 + 10\iu}{5} = -1 + 2\iu .
\]
Secara langsung: $\abs w = \sqrt{1 + 4} = \sqrt5$. Lewat kaidah hasil bagi
(\cref{prop:b1:complex:rules}~(3)): $\abs w = \frac{\abs{3 +
4\iu}}{\abs{1 - 2\iu}} = \frac5{\sqrt5} = \sqrt5$ --- jawaban yang sama,
tanpa perlu bentuk aljabar. Pelajarannya berlaku umum: \omterm{def:b1:complex:field}{modulus} dan
argumen mengalir mulus lewat hasil kali dan hasil bagi, sedangkan
bagian real dan bagian imajiner mengalir mulus lewat jumlah. Pilihlah
penyajian yang cocok dengan operasi yang sedang dihadapi, dan konversikan
hanya bila terpaksa.
\end{example}

\begin{example}[Persamaan yang melibatkan konjugat]\label{ex:b1:complex:conjeq}
Selesaikan di $\C$: $\;z + 2\conj z = 6 + 2\iu$. Persamaan yang mencampur $z$
dan $\conj z$ \emph{bukan} polinomial dalam $z$; langkah yang
andal adalah memecahnya menjadi koordinat real. Dengan $z = x + \iu y$:
\[
z + 2\conj z = 3x - \iu y ,
\]
sehingga persamaannya berbunyi $3x = 6$ dan $-y = 2$: penyelesaiannya tunggal,
yaitu $z = 2 - 2\iu$. (Periksa: $(2 - 2\iu) + 2(2 + 2\iu) = 6 +
2\iu$.) Cara lain, konjugatkan seluruh persamaannya sehingga diperoleh
$\conj z + 2z = 6 - 2\iu$ lalu selesaikan sistem linear dengan
variabel $z, \conj z$ --- jawabannya sama, dan kiat ini berguna ketika
koefisiennya kompleks. Persamaan dalam $z$ dan $\conj z$ sesungguhnya
sistem dua persamaan real; harapan ``derajat $1$, satu
penyelesaian'' aman di sini, tetapi $z\conj z = -1$ (tanpa penyelesaian) memperlihatkan
gugurnya gerak hati polinomial begitu hasil kali muncul.
\end{example}

\section{Bentuk eksponensial}

\begin{definition}[Eksponensial kompleks berargumen imajiner]\label{def:b1:complex:expi}
Untuk $\theta \in \R$ kita definisikan
\[
\eu^{\iu\theta} = \cos\theta + \iu \sin\theta .
\]
Setiap $z \neq 0$ dapat ditulis $z = r\,\eu^{\iu\theta}$ dengan $r =
\abs z > 0$; di sini $\theta$ disebut \emph{argumen}\index{argumen} dari $z$,
yang tertentu sampai penambahan kelipatan $2\pi$. Nilai yang berada di
$\intoc{-\pi}{\pi}$ disebut \emph{argumen utama}, ditulis
$\arg z$.
\end{definition}

\begin{example}[Nilai pertama, dan satu kesamaan termasyhur]\label{ex:b1:complex:eulerid}
Membaca definisi itu pada sudut-sudut istimewa:
\[
\eu^{\iu\pi/2} = \iu, \qquad
\eu^{\iu\pi} = -1, \qquad
\eu^{2\iu\pi} = 1, \qquad
\eu^{\iu\pi/4} = \frac{\sqrt2}2\,(1 + \iu) .
\]
Yang kedua, ditata ulang menjadi $\eu^{\iu\pi} + 1 = 0$, adalah kesamaan
Euler yang termasyhur, yang merangkai $\eu$, $\iu$, $\pi$, $1$ dan $0$; pada
tahap buku ini ia lebih merupakan definisi yang diurai daripada sebuah
teorema, dan isi sesungguhnya --- mengapa fungsi eksponensial analitik
pada \cref{ch:b1:functions}, yang diperluas ke argumen imajiner,
pantas menyandang nama yang sama --- diselesaikan oleh deret pangkat
pada \cref{ch:b1:series}. Sementara itu tampilan di atas
layak dihafalkan sebagai tabel konversi: ia dipakai diam-diam setiap
kali sebuah argumen dibaca dari gambar.
\end{example}

\begin{theorem}[Persamaan fungsional]\label{thm:b1:complex:funceq}
Untuk setiap $\theta, \varphi \in \R$:
\[
\eu^{\iu\theta}\, \eu^{\iu\varphi} = \eu^{\iu(\theta + \varphi)},
\qquad
\abs{\eu^{\iu\theta}} = 1,
\qquad
\conj{\eu^{\iu\theta}} = \eu^{-\iu\theta} = (\eu^{\iu\theta})^{-1}.
\]
Akibatnya $\abs{zw} = \abs z \abs w$ turut membawa argumennya:
$\arg(zw) \equiv \arg z + \arg w \pmod{2\pi}$.
\end{theorem}

\begin{proof}
Jabarkan hasil kalinya lalu pakai rumus penjumlahan sudut:
\[
(\cos\theta + \iu\sin\theta)(\cos\varphi + \iu\sin\varphi)
= (\cos\theta\cos\varphi - \sin\theta\sin\varphi)
+ \iu\,(\sin\theta\cos\varphi + \cos\theta\sin\varphi),
\]
yaitu $\cos(\theta+\varphi) + \iu\sin(\theta+\varphi)$. Modulusnya
adalah $\sqrt{\cos^2\theta + \sin^2\theta} = 1$, dan rumus \omterm{def:b1:complex:field}{konjugatnya}
merupakan keparitasan kosinus dan sinus; ia membalikkan $\eu^{\iu\theta}$ karena
$\eu^{\iu\theta}\eu^{-\iu\theta} = \eu^0 = 1$.
\end{proof}

\begin{corollary}[Rumus de Moivre]\label{cor:b1:complex:demoivre}
Untuk $\theta \in \R$ dan $n \in \Z$:
$\;(\cos\theta + \iu\sin\theta)^n = \cos(n\theta) +
\iu\sin(n\theta)$.\index{rumus de Moivre}
\end{corollary}

\begin{proof}
Untuk $n \geq 0$, berinduksilah. Kasus $n = 0$ berbunyi $1 = 1$. Dengan mengandaikan
rumus itu berlaku untuk $n$, persamaan fungsional
(\cref{thm:b1:complex:funceq}) memberikan
\[
\bigl(\eu^{\iu\theta}\bigr)^{n+1}
= \bigl(\eu^{\iu\theta}\bigr)^{n}\,\eu^{\iu\theta}
= \eu^{\iu n\theta}\,\eu^{\iu\theta}
= \eu^{\iu(n+1)\theta} ,
\]
yaitu rumus itu pada peringkat $n + 1$. Untuk $n < 0$, tulis $n = -m$
dengan $m > 0$: karena $\eu^{\iu m\theta}$ berinvers
$\eu^{-\iu m\theta}$ (\cref{thm:b1:complex:funceq} lagi),
\[
\bigl(\eu^{\iu\theta}\bigr)^{-m}
= \bigl(\eu^{\iu m\theta}\bigr)^{-1}
= \eu^{\iu(-m)\theta} . \qedhere
\]
\end{proof}

\begin{example}[Menjabarkan $\cos 3\theta$ dengan de Moivre]\label{ex:b1:complex:cos3}
Tulis $c = \cos\theta$, $s = \sin\theta$. De Moivre dan teorema
binomial memberikan
\[
\cos 3\theta + \iu \sin 3\theta = (c + \iu s)^3
= c^3 - 3cs^2 + \iu\,(3c^2 s - s^3),
\]
lalu dengan menyamakan bagian realnya, dan mensubstitusikan $s^2 = 1 - c^2$:
\[
\cos 3\theta = c^3 - 3c(1 - c^2) = 4\cos^3\theta - 3\cos\theta .
\]
Bagian imajinernya memberikan $\sin 3\theta = 3\sin\theta -
4\sin^3\theta$ cuma-cuma: satu kesamaan kompleks selalu membawa
\emph{dua} kesamaan real. Dibaca terbalik, kesamaan itu adalah
kunci klasik bagi pembagian sudut menjadi tiga: mengonstruksi
$\cos(\theta/3)$ dari $\cos\theta$ berarti menyelesaikan persamaan kubik $4x^3 -
3x = \cos\theta$, dan di situlah aljabar
\cref{ch:b1:poly} mengambil alih.
\end{example}

\begin{proposition}[Rumus Euler]\label{prop:b1:complex:euler}
\[
\cos\theta = \frac{\eu^{\iu\theta} + \eu^{-\iu\theta}}{2},
\qquad
\sin\theta = \frac{\eu^{\iu\theta} - \eu^{-\iu\theta}}{2\iu}.
\]
\end{proposition}

\begin{proof}
Jumlahkan, lalu kurangkan, $\eu^{\iu\theta} = \cos\theta +
\iu\sin\theta$ dan $\eu^{-\iu\theta} = \cos\theta - \iu\sin\theta$.
\end{proof}

\begin{method}[Trigonometri lewat eksponensial]\label{met:b1:complex:trig}
\begin{enumerate}
  \item \emph{Linearkan} $\cos^p\theta \sin^q\theta$ (ubah pangkat menjadi
        jumlah $\cos k\theta$, $\sin k\theta$): substitusikan rumus
        Euler, jabarkan dengan teorema binomial
        (\cref{thm:b1:counting:binomial}), lalu kelompokkan kembali suku yang \omterm{def:b1:complex:field}{konjugat}.
  \item \emph{Jabarkan} $\cos n\theta$ sebagai polinomial dalam $\cos\theta$:
        tulis $\cos n\theta = \Re\bigl((\cos\theta +
        \iu\sin\theta)^n\bigr)$, jabarkan, lalu ubah pangkat genap
        $\sin$ lewat $\sin^2 = 1 - \cos^2$.
  \item \emph{Jumlahkan deret trigonometri} seperti $\sum_k \cos k\theta$:
        kenali ia sebagai bagian real jumlah geometri $\sum_k
        (\eu^{\iu\theta})^k$.
  \item \emph{Pemfaktoran setengah sudut}: untuk setiap $p, q$,
        \[
        \eu^{\iu p} + \eu^{\iu q}
        = 2 \cos\tfrac{p - q}{2}\; \eu^{\iu \frac{p+q}{2}},
        \qquad
        \eu^{\iu p} - \eu^{\iu q}
        = 2\iu \sin\tfrac{p - q}{2}\; \eu^{\iu \frac{p+q}{2}} .
        \]
\end{enumerate}
\end{method}

\begin{example}[Pelinearan]\label{ex:b1:complex:linearize}
\[
\cos^3\theta
= \Bigl(\frac{\eu^{\iu\theta} + \eu^{-\iu\theta}}{2}\Bigr)^{\!3}
= \frac{\eu^{3\iu\theta} + 3\eu^{\iu\theta} + 3\eu^{-\iu\theta} +
\eu^{-3\iu\theta}}{8}
= \frac{\cos 3\theta + 3\cos\theta}{4}.
\]
Bentuk ini langsung terintegralkan --- justru itulah sebabnya pelinearan
penting pada \cref{ch:b1:integration}. Hasil kali campuran berjalan
dengan cara yang sama, hanya saja memakai kedua rumus Euler sekaligus:
\[
\sin^2\theta\cos^2\theta
= \Bigl(\frac{\sin2\theta}2\Bigr)^{\!2}
= \frac{1}{4}\cdot
\Bigl(\frac{\eu^{2\iu\theta} - \eu^{-2\iu\theta}}{2\iu}
\Bigr)^{\!2}
= \frac{2 - \eu^{4\iu\theta} - \eu^{-4\iu\theta}}{16}
= \frac{1 - \cos4\theta}{8} ,
\]
tempat jalan pintas sudut ganda pada langkah pertamanya menghemat satu
penjabaran binomial --- selalu layak dicari sebelum
mekanisasi dijalankan.
\end{example}

\begin{example}[Sebuah jumlah trigonometri binomial]\label{ex:b1:complex:binomsum}
Untuk $n \in \N$ dan $\theta \in \R$, hitunglah $S =
\sum_{k=0}^{n}\binom nk \cos k\theta$. Kenali ia sebagai bagian real
sebuah penjabaran binomial:
\[
S = \Re\sum_{k=0}^n \binom nk \bigl(\eu^{\iu\theta}\bigr)^k
= \Re\bigl(1 + \eu^{\iu\theta}\bigr)^n ,
\]
lalu faktorkan setengah sudutnya (\cref{met:b1:complex:trig}~(4)): $1 +
\eu^{\iu\theta} = 2\cos\frac\theta2\,\eu^{\iu\theta/2}$, sehingga
\[
S = \Re\Bigl(2^n\cos^n\frac\theta2\;\eu^{\iu n\theta/2}\Bigr)
= 2^n \cos^n\frac\theta2\,\cos\frac{n\theta}2 .
\]
Bagian imajinernya memberikan $\sum_k\binom nk\sin k\theta =
2^n\cos^n\frac\theta2\sin\frac{n\theta}2$ cuma-cuma. Pemeriksaan kewajaran:
$\theta = 0$ memulihkan $\sum\binom nk = 2^n$, sedangkan $\theta = \pi$
memberikan $S = 0$ untuk $n \geq 1$ (setiap faktor $\cos\frac\pi2$
lenyap), yakni jumlah baris berselang-seling pada
\cref{ex:b1:counting:binomial}. Metodenya --- ``lihatlah jumlah kosinus
itu sebagai bayangan sebuah pangkat kompleks, lalu faktorkan setengah sudutnya''
--- persis metode \cref{exo:b1:complex:6}, dengan teorema
binomial menggantikan deret geometri.
\end{example}

\section{Akar bilangan kompleks}

\begin{theorem}[Akar ke-$n$]\label{thm:b1:complex:roots}
Misalkan $a = r\,\eu^{\iu\alpha} \neq 0$ dan $n \in \N^*$. Persamaan
$z^n = a$ mempunyai tepat $n$ penyelesaian:
\[
z_k = r^{1/n}\, \eu^{\iu\left(\frac{\alpha}{n} +
\frac{2k\pi}{n}\right)},
\qquad k = 0, 1, \dots, n - 1 .
\]
\end{theorem}

\begin{proof}
Tulis $z = \rho\,\eu^{\iu\theta}$ ($\rho > 0$). Maka $z^n =
\rho^n \eu^{\iu n\theta} = r \eu^{\iu\alpha}$ jika dan hanya jika $\rho^n
= r$ (dari modulusnya) dan $n\theta \equiv \alpha \pmod{2\pi}$ (dari argumennya),
yaitu $\rho = r^{1/n}$ dan $\theta = \frac{\alpha}{n} +
\frac{2k\pi}{n}$ untuk suatu $k \in \Z$. Tinggal melihat kapan dua
bilangan bulat $k, k'$ menghasilkan bilangan yang sama: ini terjadi tepat ketika
kedua sudutnya berselisih kelipatan $2\pi$,
\[
\frac{2k\pi}n - \frac{2k'\pi}n \in 2\pi\Z
\iff \frac{k - k'}n \in \Z
\iff n \mid k - k' .
\]
Menurut pembagian Euclid, setiap $k \in \Z$ kongruen modulo $n$ dengan
tepat satu unsur $\{0, 1, \dots, n-1\}$, jadi rentang itu mendaftar
setiap penyelesaian sekali saja dan cacahannya tepat $n$. (Argumen yang
sama, dijalankan di dalam $\mathbb U_n$, menunjukkan bahwa akarnya membentuk
segi-$n$ beraturan: nilai $k$ yang berurutan memutar dengan sudut tetap
$\frac{2\pi}n$.)
\end{proof}

\begin{example}[Akar pangkat tiga dari $-27$]\label{ex:b1:complex:cuberoots}
Selesaikan $z^3 = -27$. Bentuk eksponensial ruas kanannya: $-27 =
27\,\eu^{\iu\pi}$, jadi ketiga akarnya adalah
\[
z_k = 3\,\eu^{\iu(\frac\pi3 + \frac{2k\pi}3)}, \quad k = 0, 1, 2 :
\qquad
z_0 = 3\eu^{\iu\pi/3} = \frac32 + \frac{3\sqrt3}2\,\iu,
\quad
z_1 = -3,
\quad
z_2 = \conj{z_0} .
\]
Dua pemeriksaan. Pertama, akar real $-3$ adalah yang paling kasatmata, dan
kedua akar lainnya adalah rotasinya sebesar $\pm\frac{2\pi}3$ --- setara dengan
$-3j$ dan $-3j^2$. Kedua, aljabar membenarkannya: $z^3 + 27 = (z +
3)(z^2 - 3z + 9)$, dan kuadratnya berdiskriminan $9 - 36 =
-27 < 0$ dengan akar $\frac{3 \pm 3\iu\sqrt3}2 = z_0, \conj{z_0}$.
Inti gagasannya: untuk ruas kanan yang real, akar yang tak real selalu
datang berpasangan \omterm{def:b1:complex:field}{konjugat}, sehingga gambar \omterm{def:b1:logic:sets}{himpunan} penyelesaiannya
simetris terhadap sumbu real --- pratinjau teorema pemfaktoran
real pada \cref{ch:b1:poly}.
\end{example}

\begin{example}[Ruas kanan yang tak real]\label{ex:b1:complex:fourthroots}
Selesaikan $z^4 = -8 + 8\iu\sqrt3$. Bentuk eksponensial ruas kanannya:
modulusnya $\sqrt{64 + 192} = 16$, argumennya $\theta$ dengan
$\cos\theta = -\frac12$, $\sin\theta = \frac{\sqrt3}2$, yakni
$\theta = \frac{2\pi}3$. Keempat akarnya adalah
\[
z_k = 2\,\eu^{\iu(\frac\pi6 + \frac{k\pi}2)}, \quad k = 0, 1, 2,
3 :
\qquad
z_0 = \sqrt3 + \iu,\quad z_1 = \iu z_0 = -1 + \iu\sqrt3,
\]
\[
z_2 = -z_0 = -\sqrt3 - \iu,\qquad z_3 = -\iu z_0 = 1 - \iu\sqrt3 .
\]
(Periksa: $z_0^2 = 2 + 2\iu\sqrt3$, sehingga $z_0^4 = (2 +
2\iu\sqrt3)^2 = 4 - 12 + 8\iu\sqrt3 = -8 + 8\iu\sqrt3$.) Begitu
\emph{satu} akar sudah ditemukan, ketiga akar lainnya datang cuma-cuma: semuanya
rotasi berturut-turut sebesar $\frac\pi2$, yakni hasil kalinya dengan
\omterm{def:b1:complex:unity}{akar satuan} keempat --- struktur umum di balik
\cref{thm:b1:complex:roots}, yang layak dimanfaatkan sebelum menghitung ulang
setiap akar dari nol. Tak ada simetri \omterm{def:b1:complex:field}{konjugat} kali ini:
ruas kanannya tidak real.
\end{example}

\begin{definition}[Akar satuan]\label{def:b1:complex:unity}
Sebuah \emph{akar satuan ke-$n$}\index{akar satuan} adalah
penyelesaian $z^n = 1$:
\[
\mathbb{U}_n = \bigl\{\, \omega^k : k = 0, \dots, n-1 \,\bigr\},
\qquad \omega = \eu^{2\iu\pi/n}.
\]
Mereka membentuk grup terhadap perkalian (\cref{ch:b1:structures}) dan
duduk pada titik sudut segi-$n$ beraturan yang terlukis di dalam lingkaran satuan.
\end{definition}

\begin{omfigure}
\begin{tikzpicture}[scale=1.5]
  \draw[->, thin] (-1.35,0) -- (1.45,0) node[below] {$\Re$};
  \draw[->, thin] (0,-1.35) -- (0,1.45) node[left] {$\Im$};
  \draw[thin, gray] (0,0) circle (1);
  \foreach \k in {0,...,4} {
    \fill[omDef] ({cos(72*\k)},{sin(72*\k)}) circle (0.035);
  }
  \draw[omDef, thick] (1,0) -- ({cos(72)},{sin(72)}) -- ({cos(144)},{sin(144)})
    -- ({cos(216)},{sin(216)}) -- ({cos(288)},{sin(288)}) -- cycle;
  \node[right, font=\small, omDef] at (1.02,0.12) {$1$};
  \node[above right, font=\small, omDef] at ({cos(72)},{sin(72)})
    {$\omega$};
  \node[above left, font=\small, omDef] at ({cos(144)},{sin(144)})
    {$\omega^2$};
  \node[below left, font=\small, omDef] at ({cos(216)},{sin(216)})
    {$\omega^3$};
  \node[below right, font=\small, omDef] at ({cos(288)},{sin(288)})
    {$\omega^4$};
\end{tikzpicture}

{\small \omterm{def:b1:complex:unity}{Akar satuan} kelima, $\omega = \eu^{2\iu\pi/5}$: sebuah
segi lima beraturan pada lingkaran satuan.}
\end{omfigure}

\begin{proposition}[Jumlah akar satuan]\label{prop:b1:complex:sumroots}
Untuk $n \geq 2$, \omterm{def:b1:complex:unity}{akar satuan} ke-$n$ berjumlah nol:
$\sum_{k=0}^{n-1} \omega^k = 0$.
\end{proposition}

\begin{proof}
Jumlah geometri dengan rasio $\omega \neq 1$:
$\sum_{k=0}^{n-1} \omega^k = \frac{\omega^n - 1}{\omega - 1} = 0$
karena $\omega^n = 1$.
\end{proof}

\begin{example}[Membaca bagian real dan bagian imajinernya]\label{ex:b1:complex:sumrootsparts}
Memilah $\sum_{k=0}^{n-1}\omega^k = 0$ atas bagian real dan bagian
imajiner menghasilkan dua kesamaan trigonometri cuma-cuma:
\[
\sum_{k=0}^{n-1}\cos\frac{2k\pi}n = 0,
\qquad
\sum_{k=0}^{n-1}\sin\frac{2k\pi}n = 0
\qquad (n \geq 2).
\]
Secara geometris: titik berat segi-$n$ beraturan yang terlukis di dalam
lingkaran satuan adalah pusatnya --- titik sudutnya persis saling mengimbangi.
Untuk $n = 5$ kesamaan pertama memberikan $1 + 2\cos\frac{2\pi}5 +
2\cos\frac{4\pi}5 = 0$ (dengan memasangkan $k$ dan $n - k$), yaitu titik
tolak perhitungan
$\cos\frac{2\pi}5$ pada \cref{exo:b1:complex:8}.
\end{example}

\begin{example}[Akar kuadrat dalam bentuk aljabar]\label{ex:b1:complex:sqrt}
Untuk menyelesaikan $z^2 = 3 + 4\iu$ tanpa trigonometri, tulis $z = x + \iu y$:
\[
x^2 - y^2 = 3, \qquad 2xy = 4, \qquad x^2 + y^2 = \abs{3 + 4\iu} = 5 .
\]
Menjumlahkan yang pertama dan yang terakhir: $x^2 = 4$, jadi $x = \pm 2$, lalu $y = 2/x =
\pm 1$ dengan pasangan tanda yang \emph{sama} ($xy = 2 > 0$): $z = \pm(2 +
\iu)$. Digabung dengan rumus biasa, ini menyelesaikan setiap persamaan
kuadrat berkoefisien kompleks
(\cref{exo:b1:complex:7}).
\end{example}

\section{Bilangan kompleks dan geometri bidang}

\begin{proposition}[Kamus geometri]\label{prop:b1:complex:geometry}
Identifikasikan titik $M(x, y)$ pada bidang dengan \emph{afiksnya} $z = x
+ \iu y$\index{afiks}. Untuk titik yang berbeda dua-dua $A, B, C$ berafiks $a, b,
c$:
\begin{enumerate}
  \item $\abs{b - a}$ adalah jarak $AB$;
  \item $\arg\dfrac{c - a}{b - a}$ adalah sudut antara vektor
        $\vect{AB}$ dan $\vect{AC}$ (modulo $2\pi$);
  \item $A, B, C$ segaris jika dan hanya jika $\dfrac{c - a}{b - a} \in
        \R$; garis $AB$ dan $AC$ tegak lurus jika dan hanya jika
        $\dfrac{c - a}{b - a} \in \iu\R$.
\end{enumerate}
\end{proposition}

\begin{proof}
(1) adalah definisi \omterm{def:b1:complex:field}{modulus} yang diterapkan pada $b - a$, yaitu afiks
$\vect{AB}$. (2): tulis $b - a = r\eu^{\iu\theta}$, $c - a =
s\eu^{\iu\varphi}$; maka $\frac{c-a}{b-a} = \frac sr
\eu^{\iu(\varphi - \theta)}$ berargumen $\varphi - \theta$, yaitu
sudut dari $\vect{AB}$ ke $\vect{AC}$. (3): segaris berarti sudutnya $0$
atau $\pi$, yakni argumennya di $\pi\Z$, yakni hasil baginya real;
tegak lurus berarti sudutnya $\pm\frac\pi2$, yakni hasil baginya
murni imajiner. (Hasil baginya tak nol karena $C \neq A$.)
\end{proof}

\begin{remark}[Selingan: $\C$ adalah bidang yang diberi perkalian]\label{rem:b1:complex:interlude}
Layak berhenti sejenak pada apa yang membuat bab ini mungkin sama sekali:
bidang $\R^2$ memikul penjumlahan ke segala arah, tetapi tak punya
perkalian bawaan --- dan $\C$ \emph{adalah} bidang yang
dilengkapi satu perkalian, yang di dalamnya mengalikan dengan sebuah bilangan
tetap berarti memutar sekaligus menskalakan. Struktur tunggal ini akan
diperah tiga kali lagi dalam jilid ini. Pada \cref{ch:b1:matrices}, perkalian dengan
$a + \iu b$ muncul kembali sebagai matriks $2 \times 2$ yang barisnya $(a,
-b)$ dan $(b, a)$: aritmetika kompleks adalah keluarga hasil kali matriks
yang pertama dan sepenuhnya konkret. Pada \cref{ch:b1:euclid},
rumus $\Re(\conj z\,w)$ ternyata adalah hasil kali skalar, dan
$\abs z$ adalah norma Euclid: ketaksamaan segitiga yang dibuktikan di sini
menjadi model bagi kisah Cauchy--Schwarz di sana. Dan pada
\cref{ch:b1:curves}, titik yang bergerak paling baik ditulis $t \mapsto
z(t)$, sehingga kecepatan dan percepatan menjadi turunan berkeluaran
kompleks --- gerak melingkar, misalnya, cukup ditulis $z(t) =
R\,\eu^{\iu\omega t}$. Satu perkalian yang baik, empat bab
dividennya.
\end{remark}

\begin{proposition}[Pemetaan $z \mapsto az + b$]\label{prop:b1:complex:similitude}
Misalkan $a \in \C^*$, $b \in \C$. Transformasi $f(z) = az + b$ pada
bidang:
\begin{itemize}
  \item merupakan translasi bila $a = 1$;
  \item bila tidak, ia mempunyai satu titik tetap tunggal $\zeta = \frac{b}{1-a}$,
        dan $f(z) - \zeta = a\,(z - \zeta)$: jadi $f$ adalah rotasi
        berpusat $\zeta$ dan bersudut $\arg a$, yang dikomposisikan dengan
        dilatasi berpusat $\zeta$ dan berasio $\abs a$.
\end{itemize}
Khususnya $z \mapsto \eu^{\iu\theta} z$ adalah rotasi bersudut
$\theta$ terhadap titik asal, dan $\conj z$ adalah pencerminan pada sumbu
real.
\end{proposition}

\begin{proof}
Jika $a = 1$, $f(z) = z + b$ menggeser sejauh vektor berafiks $b$. Jika
$a \neq 1$, persamaan titik tetap $z = az + b$ mempunyai penyelesaian tunggal
$\zeta = \frac{b}{1-a}$, dan lalu $f(z) - \zeta = az + b -
(a\zeta + b) = a(z - \zeta)$. Dengan menulis $a = \abs a\, \eu^{\iu\arg a}$,
perkalian dengan $a$ menskalakan jarak ke $\zeta$ sebesar $\abs a$ dan menambahkan
$\arg a$ pada sudut di $\zeta$, dan itulah komposisi yang dinyatakan tadi.
\end{proof}

\begin{example}[Menggolongkan pemetaan $z \mapsto az+b$]\label{ex:b1:complex:classify}
Ambil $f(z) = \iu z + 1$. Di sini $a = \iu \neq 1$: titik tetapnya adalah
\[
\zeta = \frac{b}{1 - a} = \frac1{1 - \iu}
= \frac{1 + \iu}{2},
\]
dan karena $\abs a = 1$, tidak ada penskalaan sama sekali: $f$ adalah rotasi
murni berpusat $\frac{1+\iu}2$ dan bersudut $\arg \iu = \frac\pi2$.
Periksa: $f(\zeta) = \iu\,\frac{1+\iu}2 + 1 = \frac{\iu - 1}2 + 1 =
\frac{1 + \iu}2 = \zeta$, lalu $f(0) = 1$, $f(1) = 1 + \iu$, $f(1 +
\iu) = \iu\,(1 + \iu) + 1 = \iu$: keempat titik $0, 1, 1+\iu, \iu$
pada persegi satuan berputar mengelilingi pusatnya $\zeta$, seperempat
putaran setiap kali --- persis yang harus dilakukan rotasi sebesar $\frac\pi2$
terhadap pusat persegi itu.
\end{example}

\begin{omfigure}
\begin{tikzpicture}[scale=1.6]
  \draw[->, thin] (-0.5,0) -- (1.6,0) node[below] {$\Re$};
  \draw[->, thin] (0,-0.35) -- (0,1.5) node[left] {$\Im$};
  \draw[gray, thin] (0,0) rectangle (1,1);
  \fill[omDef] (0.5,0.5) circle (0.03)
    node[below right, font=\small] {$\zeta$};
  \fill (0,0) circle (0.025) node[below left, font=\small] {$0$};
  \fill (1,0) circle (0.025) node[below right, font=\small] {$1$};
  \fill (1,1) circle (0.025)
    node[above right, font=\small] {$1{+}\iu$};
  \fill (0,1) circle (0.025) node[above left, font=\small] {$\iu$};
  \draw[omThm, ->, thick] (0.08,-0.12)
    .. controls (0.6,-0.28) and (1.12,0.4) .. (1.12,0.9);
  \draw[omThm, ->, thick] (1.12,1.08)
    .. controls (0.9,1.3) and (0.35,1.3) .. (0.05,1.12);
  \node[omThm, font=\small] at (0.95,-0.28) {$0 \mapsto 1$};
  \node[omThm, font=\small] at (1.62,0.75) {$1 \mapsto 1{+}\iu$};
\end{tikzpicture}

{\small \omterm{def:b1:logic:map}{Pemetaan} $f(z) = \iu z + 1$ pada
\cref{ex:b1:complex:classify}: rotasi seperempat putaran terhadap
titik tetap $\zeta = \frac{1+\iu}2$. Titik sudut persegi satuan
berputar $0 \to 1 \to 1{+}\iu \to \iu \to 0$; tak ada titik yang bergerak
sepanjang garis lurus, namun seluruh perseginya ikut terputar dengan kaku.}
\end{omfigure}

\begin{remark}[Jebakan yang lazim dengan modulus dan argumen]\label{rem:b1:complex:pitfalls}
\begin{enumerate}
  \item \emph{Tak ada ketaksamaan di $\C$.} Menulis $z \leq w$ untuk
        bilangan yang tak real tidaklah bermakna; hanya \omterm{def:b1:complex:field}{modulus}, bagian real
        dan bagian imajiner yang dapat dibandingkan.
  \item \emph{$\sqrt{\phantom z}$ disediakan bagi $\R_+$.} Setiap
        bilangan kompleks tak nol mempunyai \emph{dua} akar kuadrat dan
        tak satu pun diistimewakan: tulislah ``misalkan $\delta$ sebuah akar
        kuadrat dari $\Delta$'' (yang dihitung seperti pada
        \cref{ex:b1:complex:sqrt}), jangan pernah $\sqrt\Delta$ ---
        kaidah $\sqrt{ab} = \sqrt a\sqrt b$ sudah gugur pada
        $a = b = -1$.
  \item \emph{Argumen hidup modulo $2\pi$.} Dari
        $\eu^{\iu\alpha} = \eu^{\iu\beta}$ simpulkan $\alpha
        \equiv \beta \pmod{2\pi}$, bukan $\alpha = \beta$;
        melupakan $2k\pi$ itulah cara \omterm{def:b1:logic:sets}{himpunan} penyelesaian $z^n = a$
        kehilangan $n - 1$ dari $n$ unsurnya.
  \item \emph{$\abs{z + w}$ bukanlah $\abs z + \abs w$.} Kesamaan
        pada ketaksamaan segitiga adalah kasus segaris yang
        istimewa (\cref{prop:b1:complex:rules}~(4)); pada umumnya
        \omterm{def:b1:complex:field}{modulus} sebuah jumlah harus ditaksir, bukan dihitung.
\end{enumerate}
\end{remark}

\begin{example}\label{ex:b1:complex:equilateral}
$A, B, C$ (berafiks $a, b, c$, berbeda dua-dua) membentuk segitiga
sama sisi dengan titik sudut berurutan langsung (berlawanan arah jarum jam) jika dan
hanya jika $\frac{c - a}{b - a} = \eu^{\iu\pi/3}$: rotasi
berpusat $A$ dan bersudut $\frac\pi3$ mengirim $B$ ke $C$. Kedua orientasinya
sekaligus tertangkap oleh persamaan simetris $a^2 + b^2 + c^2 = ab
+ bc + ca$ (\cref{exo:b1:complex:10}).
\end{example}

\begin{remark}[Di mana bab ini dipakai]\label{rem:b1:complex:whereused}
Bentuk eksponensial adalah perkakas hitung yang paling banyak dipakai ulang
dalam jilid ini. \omterm{def:b1:complex:unity}{Akar satuan} menjadi contoh baku grup
siklik pada \cref{ch:b1:structures} dan menggerakkan pemfaktoran
$X^n - 1$ pada \cref{ch:b1:poly}; pengelompokan akar \omterm{def:b1:complex:field}{konjugat} di sana
menghasilkan pemfaktoran real yang dipakai pecahan parsial pada
\cref{ch:b1:fractions}. Pelinearan (\cref{met:b1:complex:trig})
adalah persiapan baku untuk mengintegralkan pangkat trigonometri pada
\cref{ch:b1:integration}, dan persamaan karakteristik pada
\cref{ch:b1:diffeq} mempunyai akar kompleks yang bagian real dan
bagian imajinernya menghasilkan penyelesaian berayun $\eu^{\lambda t}\cos\omega
t$. Kamus geometrinya kembali berbaju matriks: rotasi
dan kesebangunan menjelma menjadi matriks ortogonal pada
\cref{ch:b1:matrices,ch:b1:euclid}, dan kurva terparameter pada
\cref{ch:b1:curves} sering paling baik ditulis sebagai \omterm{def:b1:logic:map}{pemetaan} $t \mapsto z(t)
\in \C$.
\end{remark}

\section{Latihan}

\begin{exercise}[$\star$]\label{exo:b1:complex:1}
Nyatakan dalam bentuk eksponensial: $1 + \iu$; $\;\sqrt 3 - \iu$; $\;-5$;
$\;\dfrac{1 + \iu}{\sqrt 3 - \iu}$. Lalu simpulkan
$\cos\frac{5\pi}{12}$ dan $\sin\frac{5\pi}{12}$.
\end{exercise}

\begin{exercise}[$\star$]\label{exo:b1:complex:2}
Hitung $(1 + \iu)^{20}$. Untuk $n \in \N$ yang mana $(1 + \iu)^n$ merupakan
bilangan real?
\end{exercise}

\begin{exercise}[$\star$]\label{exo:b1:complex:3}
Gambarkan secara geometris \omterm{def:b1:logic:sets}{himpunan} $z \in \C$ yang memenuhi:
$\abs{z - 2} = \abs{z + \iu}$; $\;\abs{z - 1} = 2$;
$\;\dfrac{z - 1}{z + 1} \in \iu\R$ (untuk $z \neq -1$).
\end{exercise}

\begin{exercise}[$\star$]\label{exo:b1:complex:4}
Linearkan $\sin^4\theta$, lalu jabarkan $\cos 4\theta$ sebagai polinomial dalam
$\cos\theta$.
\end{exercise}

\begin{exercise}[$\star$]\label{exo:b1:complex:5}
Selesaikan $z^3 = 8\iu$ lalu tempatkan penyelesaiannya pada sebuah gambar. Selesaikan
$z^4 = -4$ lalu faktorkan $X^4 + 4$ menjadi dua polinomial kuadrat real.
\end{exercise}

\begin{exercise}[$\star\star$]\label{exo:b1:complex:6}
Untuk $\theta \in \R$ dengan $\eu^{\iu\theta} \neq 1$ dan $n \in \N$,
hitunglah
\[
C_n = \sum_{k=0}^{n} \cos k\theta
\qquad\text{dan}\qquad
S_n = \sum_{k=0}^{n} \sin k\theta
\]
dengan menjumlahkan deret geometri $\sum_k \eu^{\iu k\theta}$ lalu memakai
pemfaktoran setengah sudut.
\end{exercise}

\begin{exercise}[$\star\star$]\label{exo:b1:complex:7}
Selesaikan di $\C$: $z^2 - (3 + 4\iu) z + (-1 + 5\iu) = 0$. \emph{(Hitung
diskriminannya, lalu tarik akar kuadratnya seperti pada
\cref{ex:b1:complex:sqrt}.)}
\end{exercise}

\begin{exercise}[$\star\star$]\label{exo:b1:complex:8}
Misalkan $\omega = \eu^{2\iu\pi/5}$.
\begin{enumerate}
  \item Berikan alasan bahwa $1 + \omega + \omega^2 + \omega^3 + \omega^4 = 0$.
  \item Tulis $u = \omega + \omega^4$ dan $v = \omega^2 + \omega^3$.
        Hitung $u + v$ dan $uv$, lalu simpulkan bahwa $u$ dan $v$ adalah
        akar $X^2 + X - 1$.
  \item Simpulkan bahwa $\cos\frac{2\pi}{5} = \frac{\sqrt 5 - 1}{4}$.
\end{enumerate}
\end{exercise}

\begin{exercise}[$\star\star$]\label{exo:b1:complex:9}
Buktikan bahwa untuk setiap $z, w \in \C$ berlaku (kesamaan jajargenjang):
\[
\abs{z + w}^2 + \abs{z - w}^2 = 2\abs z^2 + 2\abs w^2 ,
\]
lalu tafsirkan secara geometris pada jajargenjang bertitik sudut $0,
z, w, z + w$.
\end{exercise}

\begin{exercise}[$\star\star\star$]\label{exo:b1:complex:10}
Buktikan bahwa tiga titik yang berbeda dua-dua dan berafiks $a, b, c$ membentuk
segitiga sama sisi (dengan orientasi apa pun) jika dan hanya jika
\[
a^2 + b^2 + c^2 = ab + bc + ca .
\]
\emph{Petunjuk: kasus langsungnya $\frac{c-a}{b-a} = -j^2$ dan
kasus tak langsungnya $\frac{c-a}{b-a} = -j$, dengan $j = \eu^{2\iu\pi/3}$
memenuhi $j^2 + j + 1 = 0$; atau faktorkan $a + jb + j^2c$ dan $a + j^2 b
+ jc$.}
\end{exercise}

\begin{exercise}[$\star\star\star$]\label{exo:b1:complex:11}
Untuk $n \in \N^*$, hitunglah $P = \prod_{k=1}^{n-1}
\bigl(1 - \omega^k\bigr)$ dengan $\omega = \eu^{2\iu\pi/n}$.
\emph{Petunjuk: $X^n - 1 = \prod_{k=0}^{n-1} (X - \omega^k)$; bagilah dengan
$X - 1$ lalu hitung nilainya di $X = 1$.} Lalu simpulkan
$\prod_{k=1}^{n-1} \sin\frac{k\pi}{n} = \dfrac{n}{2^{n-1}}$.
\end{exercise}

\begin{exercise}[$\star\star$]\label{exo:b1:complex:12}
Misalkan $n \geq 2$. Selesaikan persamaan $(z + 1)^n = (z - 1)^n$ di $\C$:
tunjukkan bahwa ia mempunyai tepat $n - 1$ penyelesaian, semuanya murni imajiner,
yaitu
\[
z_k = -\iu\,\frac{\cos\frac{k\pi}{n}}{\sin\frac{k\pi}{n}},
\qquad k = 1, \dots, n - 1 .
\]
\emph{Petunjuk: $z = 1$ bukan penyelesaian, jadi bagilah lalu pakai \omterm{def:b1:complex:unity}{akar
satuan}; setelah itu terapkan pemfaktoran setengah sudut pada
\cref{met:b1:complex:trig}.}
\end{exercise}

\section{Soal: Teorema Napoleon}

\begin{problem}\label{pb:b1:complex:1}
Dirikan segitiga sama sisi ke arah luar pada setiap sisi sebuah
segitiga \emph{sembarang}: ketiga pusat segitiga itu
selalu membentuk segitiga sama sisi. Inilah teorema Napoleon
\index{teorema Napoleon} --- \omterm{def:b1:logic:statement}{pernyataan} yang tak punya alasan kasatmata untuk
menjadi benar, yang dibuktikan aljabar $j = \eu^{2\iu\pi/3}$ dalam tiga
baris perhitungan. Soal ini membangun perkakas lengkapnya
(rotasi, kesebangunan langsung, kriteria $j$ pada
\cref{exo:b1:complex:10}), membuktikan teorema Napoleon yang dalam maupun
yang luar, menentukan letak kasus merosotnya, lalu menutup dengan permata kedua
dari mazhab yang sama: teorema van Schooten tentang segitiga sama sisi
dan ketaksamaan Ptolemeus. Di sepanjang soal ini, titik bidang
diidentifikasikan dengan afiksnya.

\medskip
\noindent\textbf{Bagian I --- Bilangan $j$ dan rotasi.}
\begin{enumerate}
  \item Hitung bentuk aljabar $j$, lalu $j^2$, $j^3$, $1 + j
        + j^2$, $\conj j$ dan $j^{-1}$. Tempatkan $1$, $j$, $j^2$ pada
        sketsa lingkaran satuan.
  \item Tunjukkan bahwa rotasi berpusat $a$ dan bersudut $\theta$ adalah
        $r(z) = a + \eu^{\iu\theta}(z - a)$, dan bahwa
        komposisi dua rotasi, yang bersudut $\theta$ dan
        $\theta'$, merupakan rotasi bersudut $\theta + \theta'$ bila
        $\theta + \theta' \notin 2\pi\Z$, dan merupakan translasi
        bila tidak. (Pakai \cref{prop:b1:complex:similitude}.)
  \item Misalkan $b \neq c$. Tunjukkan bahwa tepat ada dua titik $p$
        yang membuat $(b, c, p)$ sama sisi, yaitu $p = b +
        \eu^{\pm\iu\pi/3}(c - b)$. Untuk segitiga \emph{langsung}
        $(a, b, c)$, periksa pada contoh $a = 0$, $b = 1$, $c =
        \iu$ bahwa pilihan $\eu^{-\iu\pi/3}$ adalah yang terletak di
        seberang garis $BC$ dari $a$ --- yaitu puncak
        \emph{arah luarnya}.
  \item Dari penyelesaian \cref{exo:b1:complex:10}, $(a, b, c)$
        sama sisi langsung jika dan hanya jika $a + jb + j^2c = 0$. Buktikan kedua
        kesamaan rotasi
        \[
        b + jc + j^2a = j^2\,(a + jb + j^2c),
        \qquad
        c + ja + j^2b = j\,(a + jb + j^2c),
        \]
        lalu simpulkan bahwa kriteria itu tak berubah di bawah
        \omterm{def:b1:counting:objects}{permutasi} siklik $(a, b, c)$.
  \item Tunjukkan bahwa jika $a + jb + j^2c = 0$ dan dua di antara ketiga titiknya
        berimpit, maka ketiganya berimpit. (Jadi kriteria itu persis
        mencirikan: segitiga sama sisi langsung, atau sebuah titik
        tunggal.)
\end{enumerate}

\medskip
\noindent\textbf{Bagian II --- Kesebangunan langsung.}
\begin{enumerate}[resume]
  \item Tunjukkan bahwa \omterm{def:b1:logic:map}{pemetaan} $f(z) = az + b$ dengan $a \in \C^*$
        (\emph{kesebangunan langsung}) tertutup terhadap komposisi
        dan pembalikan: komposisi dua di antaranya, dan
        invers salah satunya, kembali berbentuk sama. (Dalam
        bahasa \cref{ch:b1:structures}: mereka membentuk grup.)
  \item Diberikan $z_1 \neq z_2$ dan $w_1 \neq w_2$, tunjukkan bahwa ada
        \emph{tepat satu} kesebangunan langsung $f$ dengan $f(z_1) = w_1$ dan
        $f(z_2) = w_2$, lalu berikan $a$ dan $b$ secara eksplisit.
  \item Tunjukkan bahwa kesebangunan langsung $f(z) = az + b$ mengalikan
        semua jarak dengan $\abs a$ dan mengawetkan \emph{bentuk}
        $\frac{c - a'}{b' - a'}$ sebarang segitiga $(a', b', c')$
        \emph{(pembilang dan penyebutnya sama-sama dikalikan
        $a$)}. Simpulkan bahwa dua segitiga sebangun langsung
        tepat ketika bentuknya sama.
  \item Tentukan selengkapnya kesebangunan langsung dengan $f(0) = 1$
        dan $f(1) = \iu$: berikan $a$, $b$, titik tetapnya, rasionya
        dan sudutnya.
  \item Misalkan $\alpha + \beta = 1$. Tunjukkan bahwa ``titik berbobot''
        $g(a', b') = \alpha a' + \beta b'$ berkomutasi dengan setiap
        kesebangunan langsung: $f(\alpha a' + \beta b') = \alpha f(a')
        + \beta f(b')$. Simpulkan bahwa titik berat, titik tengah, dan
        pusat Napoleon di bawah semuanya terangkut oleh kesebangunan
        --- yaitu izin aljabar di balik setiap argumen ``tanpa mengurangi
        keumuman, letakkan lingkaran luarnya pada lingkaran satuan''.
\end{enumerate}

\medskip
\noindent\textbf{Bagian III --- Teorema Napoleon.} Misalkan $(a, b, c)$
sebuah segitiga langsung. Pada setiap sisinya dirikan segitiga sama sisi
ke arah luar (pertanyaan 3) lalu misalkan $n_a$, $n_b$, $n_c$ pusat
(titik berat) segitiga yang didirikan pada $[b, c]$, $[c, a]$,
$[a, b]$ berturut-turut.
\begin{enumerate}[resume]
  \item Tunjukkan bahwa
        \[
        n_a = \alpha b + \beta c, \qquad
        n_b = \alpha c + \beta a, \qquad
        n_c = \alpha a + \beta b,
        \qquad\text{dengan}\quad
        \alpha = \frac{3 + \iu\sqrt3}{6},\ \beta = \conj\alpha .
        \]
  \item Periksa $\alpha + \beta = 1$, lalu simpulkan bahwa segitiga
        $(n_a, n_b, n_c)$ mempunyai \emph{titik berat yang sama} dengan
        $(a, b, c)$.
  \item Dengan memakai kesamaan pada pertanyaan 4, tunjukkan bahwa
        \[
        n_a + j\,n_b + j^2 n_c
        = (\alpha j^2 + \beta j)\,(a + jb + j^2c) .
        \]
  \item Hitung $\alpha j + \beta$ lalu simpulkan: $\alpha j^2 +
        \beta j = j(\alpha j + \beta) = 0$, sehingga $n_a + j n_b + j^2
        n_c = 0$ untuk \emph{setiap} segitiga: jadi segitiga Napoleon
        luarnya sama sisi langsung (atau sebuah titik). Teorema Napoleon
        pun terbukti.
  \item Dirikan segitiga sama sisinya \emph{ke arah dalam} sebagai gantinya
        (pilihan $\eu^{+\iu\pi/3}$ pada pertanyaan 3) lalu misalkan $m_a, m_b,
        m_c$ pusatnya. Tunjukkan $m_a = \beta b + \alpha c$ (dan
        seterusnya secara siklik), lalu buktikan $m_a + j^2 m_b + j\,m_c = 0$:
        jadi segitiga Napoleon dalamnya pun sama sisi, dengan orientasi
        yang berlawanan.
  \item Tentukan letak kemerosotannya: tunjukkan bahwa $n_a = n_b = n_c$ terjadi
        tepat ketika $a + j^2 b + jc = 0$, yakni ketika $(a, b, c)$
        merupakan segitiga sama sisi \emph{tak langsung}. \emph{(Dua
        titik di antara $n_a, n_b, n_c$ berimpit jika dan hanya jika ketiganya berimpit, menurut
        pertanyaan 5; lalu pakai $n_a + j^2 n_b + j n_c = j(\alpha +
        \beta j)(a + j^2b + jc)$ dan periksa bahwa $\alpha + \beta j \neq
        0$.)}
  \item Jalankan seluruh perhitungannya pada segitiga $a = 0$,
        $b = 1$, $c = \iu$: berikan $n_a, n_b, n_c$ secara eksak, lalu
        buktikan lewat perhitungan langsung atas ketiga kuadrat panjang
        sisinya bahwa segitiga itu sama sisi, dengan kuadrat sisi
        $\frac{2 + \sqrt3}{3}$.
\end{enumerate}

\medskip
\noindent\textbf{Bagian IV --- Ptolemeus dan van Schooten.}
\begin{enumerate}[resume]
  \item Buktikan kesamaan berikut, yang berlaku untuk setiap $a, b, c, d$ kompleks:
        \[
        (a - b)(c - d) + (a - d)(b - c) = (a - c)(b - d) .
        \]
  \item Simpulkan \emph{ketaksamaan Ptolemeus}\index{ketaksamaan
        Ptolemeus}: untuk sebarang empat titik $A, B, C, D$,
        \[
        AC \cdot BD \;\leq\; AB \cdot CD + AD \cdot BC ,
        \]
        dengan kesamaan jika dan hanya jika $(a-b)(c-d)$ dan
        $(a-d)(b-c)$ terletak pada satu sinar garis yang sama dari $0$.
  \item Tunjukkan bahwa untuk $\theta, \varphi \in \R$ berlaku
        $\abs{\eu^{\iu\theta} - \eu^{\iu\varphi}} =
        2\,\abs{\sin\frac{\theta - \varphi}2}$.
  \item (Teorema van Schooten\index{teorema van Schooten}) Misalkan
        $(a, b, c) = (1, j, j^2)$ --- menurut Bagian II ini tidak mengurangi
        keumuman di antara segitiga sama sisi langsung --- lalu misalkan
        $p = \eu^{\iu\theta}$ dengan $\theta \in
        \intoo{2\pi/3}{4\pi/3}$, yaitu titik pada lingkaran luarnya di
        busur $BC$ yang tidak memuat $A$. Buktikan
        \[
        PA = PB + PC .
        \]
        \emph{(Nyatakan ketiga jaraknya dengan pertanyaan 20 lalu pakai
        rumus jumlah-ke-hasil-kali.)}
  \item Periksa pada $p = -1$ bahwa ini persis kasus kesamaan
        pada ketaksamaan Ptolemeus untuk urutan siklik $A, B, P, C$:
        hitung $(a - b)(p - c)$ dan $(a - c)(b - p)$ lalu periksa
        bahwa rasionya bilangan real positif.
\end{enumerate}

\medskip
\noindent\textbf{Bagian V --- Sintesis.}
\begin{enumerate}[resume]
  \item Hitung segitiga Napoleon luar dari segitiga sama sisi
        $(1, j, j^2)$ itu sendiri, lalu gambarkan hasilnya
        secara geometris.
  \item Di mana persisnya soal ini memakai: (i) perkalian sebagai
        rotasi; (ii) struktur grup kesebangunan beserta
        bentuk yang invarian; (iii) pemfaktoran setengah sudut pada
        \cref{met:b1:complex:trig}? Satu kalimat untuk masing-masing.
  \item Dalam satu paragraf pendek, rumuskan moral soal ini:
        apa yang dibeli kamus antara geometri bidang dan
        aljabar $\C$, apa harganya, dan langkah bukti mana di antara
        keduanya --- kesamaan $\alpha j + \beta = 0$ atau
        gambar klasiknya --- yang menurut Anda lebih baik
        \emph{menjelaskan} teorema Napoleon? Sebutkan satu tempat yang
        kamus itu akan muncul kembali di dalamnya dalam bentuk matriks
        di bagian akhir jilid ini.

\end{enumerate}
\end{problem}
